\documentclass[11pt,letterpaper]{article}
\usepackage[T1]{fontenc}
\usepackage{lmodern,microtype}
\usepackage[margin=1in]{geometry}
\usepackage{amsmath,amssymb,amsthm}
\usepackage[hidelinks]{hyperref}

\newtheorem{theorem}{Theorem}
\newtheorem{lemma}{Lemma}[section]
\newcommand{\E}{\mathbb E}
\newcommand{\Z}{\mathbb Z}
\newcommand{\one}{\mathbf 1}
\newcommand{\tr}{\operatorname{tr}}
\newcommand{\Var}{\operatorname{Var}}
\newcommand{\aut}{\operatorname{aut}}
\newcommand{\abs}[1]{\lvert #1\rvert}
\newcommand{\norm}[1]{\lVert #1\rVert}
\newcommand{\op}{\mathrm{op}}
\newcommand{\doi}[1]{\href{https://doi.org/#1}{\nolinkurl{#1}}}

\title{The sharp Lov\'asz theta constant for two random graph models}
\author{Shanmu Jin\thanks{Department of Neurosurgery, Peking Union Medical College
  Hospital, No.~1 Shuaifuyuan, Dongcheng District, Beijing
  100730, China.
  \href{mailto:jinshanmu1129@pku.edu.cn}{\nolinkurl{jinshanmu1129@pku.edu.cn}}.
  \href{https://orcid.org/0009-0009-5346-753X}{ORCID 0009-0009-5346-753X}.}}
\date{}
\hypersetup{pdftitle={The sharp Lov\'asz theta constant for two random graph models},
  pdfauthor={Shanmu Jin},
  pdfsubject={Lov\'asz theta asymptotics for independent-edge and circulant random graphs}}

\begin{document}
\maketitle

\begin{abstract}
The expected Lov\'asz theta number is asymptotic to \(\sqrt n\) both for the
Erd\H{o}s--R\'enyi graph \(G(n,1/2)\) and for the random circulant graph in
which each cyclic distance is included independently with probability
\(1/2\). For independent edges we prove a sharp norm bound for the graph
matrices in a positive semidefinite certificate. For circulant graphs we
control a smoothed Fourier program using Gaussian approximations indexed
by rooted trees. The circulant proof applies to every order, including
even orders with a single antipodal distance.
\end{abstract}

\section{Introduction}

Lov\'asz introduced the theta number as an upper bound on the Shannon
capacity of a graph~\cite{Lovasz}. For a simple graph \(G\) on \(n\)
vertices, we use the semidefinite formulation
\begin{equation}\label{eq:theta}
 \vartheta(G)=\max\{\langle J,X\rangle : X\succeq0,\ \tr X=1,\
        X_{ij}=0\text{ for }\{i,j\}\in E(G)\},
\end{equation}
where \(X\) is real symmetric and \(J\) is the all-ones matrix.
Juh\'asz established the order \(\sqrt n\) for \(G(n,1/2)\)~\cite{Juhasz}.
Potechin and Xu obtained the leading constant one by constructing a
certificate from graph matrices~\cite{PX}. For random circulants,
Bandeira et al. used a Fourier linear program and
conjectured the same constant~\cite{BBDFKK}. Chen et al. proved the
expectation limit over odd orders using smoothed Fourier equations and
rooted-tree fields~\cite{CHHWY}.

Section~\ref{sec:independent} gives a direct trace count and product rule
for the graph matrices in the independent-edge certificate. In
Section~\ref{sec:circulant}, a rank count for the leaf sums of a tree gives
Gaussian comparison over any cyclic group, including the singleton
distance class at even order.

\begin{theorem}\label{thm:main}
Let \(G_n\) have either of the following distributions.
\begin{enumerate}
\item \(G_n=G(n,1/2)\), with independent edges.
\item \(G_n\) is the graph on \(\Z_n=\Z/n\Z\) in which \(i\) and \(j\)
are adjacent if \(i-j\) lies in a random connection set. Each nonzero
inversion class \(\{s,-s\}\) is included independently with probability
\(1/2\).
\end{enumerate}
Then \(\E\vartheta(G_n)/\sqrt n\to 1\) as \(n\to\infty\) through all
positive integers.
\end{theorem}

The lower bound is common to the two models. Strict feasibility at
\(I/n\) in \eqref{eq:theta} gives a dual optimum
\(Q=\vartheta(G)I-J+Y\succeq0\), with \(Y\) supported on the edges of
\(G\). The matrix \((Q+J)/(n\vartheta(G))\) is feasible for the complement
\(\overline G\). Since \(\langle J,Q\rangle\ge0\), it follows that
\(\vartheta(G)\vartheta(\overline G)\ge n\)~\cite{Lovasz}. Both distributions are
invariant under complementation, so \(\E\vartheta(G_n)\ge\sqrt n\) by
the arithmetic--geometric mean inequality. We will also use the bounds
\(1\le\vartheta(G)\le n\), which follow from \(I/n\) and \(J\preceq nI\).

\section{Independent edges}\label{sec:independent}

Let \(A\) be zero on the diagonal, \(+1\) on edges, and \(-1\) on
nonedges. Its entries above the diagonal are independent fair signs.
A matrix \(W\succeq0\) with constant diagonal \(b\) and entries
\(-a/\sqrt n\) on off-diagonal nonedges, where \(a>0\), gives
\begin{equation}\label{eq:witness}
 \vartheta(G_n)\le 1+(b/a)\sqrt n.
\end{equation}
Indeed, \(B=(1+b\sqrt n/a)I-(\sqrt n/a)W\) agrees with \(J\) on every
entry that can contribute to the objective in \eqref{eq:theta}, and
\(B\preceq(1+b\sqrt n/a)I\). We construct such a matrix with \(b/a\)
close to one.

\subsection{Graph matrices}

A shape is a finite simple graph with ordered left and right boundary
vertices, which may coincide. Its graph matrix has entries
\[
 (M_\alpha)_{ij}=
 \sum_{\substack{\phi:V(\alpha)\hookrightarrow\{1,\ldots,n\}\\
                 \phi(\mathrm{left})=i,\ \phi(\mathrm{right})=j}}
          \prod_{\{x,y\}\in E(\alpha)} A_{\phi(x),\phi(y)}.
\]
Shapes are taken up to isomorphism preserving the ordered boundaries.
The basic shapes are the single edge and the graphs formed from a cycle
by adding noncrossing chords, with the two boundaries adjacent on the
cycle~\cite{PX}. The cycle is the unique spanning cycle. Indeed, the
endpoints of a chord separate the interiors of its two arcs, so a
spanning path between these endpoints cannot traverse both interiors.
In particular, an automorphism fixing both boundaries fixes the whole
shape. A chain is formed by identifying the right boundary of each
basic shape with the left boundary of its successor. A single vertex
is a chain of length zero. Write \(\widehat M_\alpha\) for
\(n^{-(\abs{V(\alpha)}-1)/2}M_\alpha\) for a chain, and for
\(n^{-\abs{V(\alpha)}/2}M_\alpha\) for two disjoint chains with one
outer boundary on each.

We will use a consequence of the spanning cycle. When the edges of a
basic shape are split into two nonempty sets, at least two vertices
meet both sets. If there are exactly two such vertices, each set of
edges gives a connected graph on its incident vertices. Indeed,
deleting a vertex from the spanning cycle leaves a spanning path, so
every component of either edge set must meet the other set in at least
two vertices.

\begin{lemma}\label{lem:rough}
For a fixed shape \(\gamma\), let \(v\) be its vertex count and \(z\)
the number of isolated vertices outside its boundaries. Let \(\iota=1\) if
its boundaries are connected or coincide, and \(\iota=0\) otherwise. For
any fixed \(m>0\), with probability \(1-O(n^{-m})\),
\begin{equation}\label{eq:rough}
 \norm{M_\gamma}_{\op}\le C(\log n)^C n^{(v+z-\iota)/2},
\end{equation}
where \(C\) depends only on \(\gamma\) and \(m\).
\end{lemma}

\begin{proof}
This is the graph-matrix separator bound with one vertex on each
boundary~\cite{AMP}. Color the \(n\) labels independently and uniformly
by the vertices of \(\gamma\), and restrict each vertex to labels of
its color. The resulting matrix \(M^{\mathrm{col}}\) satisfies
\(M_\gamma=v^v\E_{\mathrm{col}} M^{\mathrm{col}}\). Convexity bounds the
\(2\ell\)-th norm moment of \(M_\gamma\) by \(v^{2\ell v}\) times the
average of \(\E\tr[(M^{\mathrm{col}}(M^{\mathrm{col}})^{\mathsf T})^\ell]\).

For a fixed coloring, every edge in a contributing trace term must
repeat in copies of the same shape edge. Thus each nonisolated vertex
role uses at most \(\ell\) different labels across the \(2\ell\) copies.
Each boundary role also uses at most \(\ell\) labels because the trace
identifies it in pairs. Isolated interior vertices use at most \(2\ell\)
labels per role. If \(\iota=1\), fix a path of length \(j\) joining the
boundaries, taking \(j=0\) for coincident boundaries. Its copies form a
closed walk with at most \(\ell j\) distinct edges, hence at most \(\ell j+1\)
labels altogether. The total label count is therefore at most
\(\ell(v+z-\iota)+\iota\). There are at most \((2\ell v)^{2\ell v}\)
equality patterns of labels. Taking \(\ell=\lceil\log n\rceil\) and
applying Markov's inequality gives \eqref{eq:rough}.
\end{proof}

\begin{lemma}\label{lem:sharp}
Let \(M_c=\sum_\alpha c_\alpha\widehat M_\alpha\) be a symmetric
combination of finitely many basic shapes with fixed real coefficients.
For every \(\epsilon>0\) and fixed \(m>0\), with probability
\(1-o(n^{-m})\),
\begin{equation}\label{eq:sharp}
 \norm{M_c}_{\op}\le(2+\epsilon)\Bigl(\sum_\alpha c_\alpha^2\Bigr)^{1/2}.
\end{equation}
\end{lemma}

\begin{proof}
By homogeneity take \(\sum_\alpha c_\alpha^2=1\), and let \(v_\ast\) be the
largest shape size. Expand \(\E\tr M_c^\ell\) for even \(\ell\).
A term specifies a closed sequence of \(\ell\) labeled basic shapes,
called blocks in this count. Consecutive blocks share the right label
of one and the left label of the next. Only terms in which every edge
occurs an even number of times contribute. We count their equality
patterns before choosing the actual labels.

Let \(m_i\) count the labels in block \(i\) that occur both earlier and
later in the sequence. The edge between the two boundaries must
repeat, so \(m_i\ge1\) for \(1<i<\ell\). Put
\(t=\sum_{1<i<\ell}(m_i-1)\), and let \(v_i\) be the size of block
\(i\). Each label occurs at least twice. Counting half a label at
its first occurrence and half at its last gives
\[
 \#\{\text{labels}\}=1+\tfrac12\sum_{i=1}^\ell(v_i-1)-t/2.
\]
Including the normalization of each matrix, a pattern contributes at
most \(n^{1-t/2}\) times its absolute coefficient weight.

There are at most \(t\) blocks with \(m_i\ge2\). Every other block is
of one of two types. In a forward block only the left label can have
occurred earlier. In a return block only the right label can occur
later. The first block is forward and the last is a return. For an
interior block with \(m_i=1\), its repeated boundary edge forces that
unique label to be a boundary. If it is the left boundary, the right
label occurs only later. An edge between an earlier-only and a
later-only label would occur once. Deleting the left boundary leaves
a connected graph, so every other label occurs only later. Reversing
the order proves the return case.

For a forward block \(B_i\), let \(B_j\) be the first later block that
shares an edge with it. Suppose \(m_j<2\), so \(B_j\) is a return block.
Indeed, the shared edge gives \(B_j\) two labels that have already occurred.
If an edge of \(B_i\) were absent from \(B_j\), splitting the edges of
\(B_i\) by membership in \(B_j\) would give two mixed vertices. Each occurs
in \(B_j\) and has a later occurrence, because the edge absent from
\(B_j\) has yet to repeat. This forces \(m_j\ge2\). Hence
\(E(B_i)\subseteq E(B_j)\). If the inclusion is strict, let \(B_k\) be the
last block before \(B_j\) sharing an edge with it. Then \(i\le k<j\).
Some edge of \(B_j\) lies outside \(B_k\). For \(k=i\) this follows from
strict inclusion, and for \(k>i\) it follows because \(B_k\) misses every
edge of \(B_i\) by the choice of \(j\). Split the edges of \(B_j\) by
membership in \(B_k\). Each mixed vertex occurs in \(B_j\) and before
\(B_k\), since an edge of \(B_j\) outside \(B_k\) must repeat before
\(B_k\). Thus \(m_k\ge2\), which forces \(i<k<j\) because \(B_i\) is forward.

Pair \(B_i\) and \(B_j\) when \(m_j<2\) and \(E(B_i)=E(B_j)\). Their edges
have no earlier occurrences by the forward rule, none in between by
the choice of \(j\), and none later by the return rule.
The paired blocks also have the same labels, since every vertex is
incident to an edge.
Reversing the argument for return blocks shows that
every unpaired block is at most two edge-sharing steps from some
block with \(m_i\ge2\), with any intermediate block of forward or
return type. A labeled edge has at most one occurrence in a forward
block and one in a return block, since those are its first and last
occurrences. As each block has at most \(\binom{v_\ast}{2}\) edges, the
number of unpaired blocks is \(O_{v_\ast}(t)\).

Mark the unpaired blocks and both members of each pair whose forward
block has an internal label that occurs in another block before its
return. Take the first intervening occurrence of such a label. A
paired forward block could contain it only at its left boundary, so
its predecessor would have it at its right boundary. The original
forward block has it internally, and a later predecessor would
contradict the choice of the first occurrence. A paired return block
from a different pair would require its own forward partner to have
used the label before the original forward block or between that
block and the first intervening occurrence. Thus the first
intervening occurrence is in an unpaired block.
Each such block accounts for at most \(v_\ast\)
additional pairs, since an internal label is new in its forward
block. Hence the number \(r\) of marked positions is \(O_{v_\ast}(t)\).

Fix these positions and record the equality pattern in sequence.
Give each newly exposed label its name in order of exposure. A marked
block exposes all its labels in a fixed vertex order of its shape.
An unmarked block exposes only its boundaries. At a marked step,
specify the shape and record for each vertex a previous name or a
new label, using at most \((v_\ast\ell+1)^{v_\ast}\) choices beyond the
shape. At an unmarked forward step, store a waiting pair with its two ordered
endpoints, giving the right endpoint a new name. Its internal labels
stay private to that pair until the return.

At an unmarked return, the preceding block supplies the starting
label, which must therefore be an endpoint of the pair. The only
previous block to use a private label is its own forward block, and
that block ends at an endpoint. The pair is unique among the waiting
pairs at the start, because the starting label cannot recur after
the return.
The destination is either the pair's other endpoint or the other
cycle neighbor of the start, if present. If another waiting pair
contains the other endpoint, that endpoint must be the destination
because it occurs in a later return. Otherwise neither destination
can start an immediately following unmarked return. Thus one bit at
each step records the next unmarked type when needed, or the
destination at a return when the next type is forced. The first
position, if unmarked, is forward. After a return, only its
destination can recur among the labels of the pair. In particular,
if the destination is internal to the forward block, give it a new
name at the return.

The exposed names, the pairing, and each change of boundaries are
therefore determined by the marked shapes and the code alone. Number
the vertices of each forward cycle from its left boundary in the
direction of its right boundary. For a fixed size and change of
boundaries, reversing or rotating these cyclic positions gives the
return order. This permutation is fixed by the code and invertible.
Thus the return shape is \(\rho(\alpha)\) for the forward shape
\(\alpha\), with \(\rho\) injective on the allowed shapes. The sum of
the pair's absolute weights is at most
\(\sum_\alpha\abs{c_\alpha c_{\rho(\alpha)}}\le1\) by Cauchy--Schwarz,
where coefficients outside the given family are zero. Once the
forward shapes are chosen, assign the private labels to their
respective positions in each pair to recover the full equality
pattern. For \(t=0\) the total weight is at most
\(2^\ell\). For \(t\ge1\), choosing and describing \(O_{v_\ast}(t)\) marked
positions, and summing their shape weights, gives the bound
\(2^\ell(v_\ast\ell+1)^{C t}\), with \(C\) fixed. Taking
\(\ell=2\lceil(\log n)^2\rceil\) gives
\[
 \E\tr M_c^\ell\le n\,2^\ell
    \left(1+\sum_{t\ge1}\left[\frac{(v_\ast\ell+1)^C}{\sqrt n}\right]^t\right)
    =n\,2^\ell(1+o(1)).
\]
Markov's inequality proves \eqref{eq:sharp}.
\end{proof}

\subsection{The chain product}

\begin{lemma}\label{lem:product}
Multiply the normalized graph matrix of a fixed chain or a fixed pair
of disjoint chains on the right by that of a fixed basic shape
\(\sigma\). Up to an error \(o(1)\) in operator norm with probability
\(1-O(n^{-2})\), the product is the sum of the following terms, each
with coefficient one for each identification.
\begin{enumerate}
\item Append \(\sigma\) at the right boundary, with all other vertices
      distinct.
\item Identify \(\sigma\) edge for edge with the last block of the right
      chain. If the new right boundary is that block's left endpoint,
      delete the block. For the other possible right boundary, a
      single chain becomes the shortened chain and a disjoint isolated
      right boundary. For a pair of chains this last term occurs when
      the right chain had length one, and replaces it by an isolated
      right boundary.
\end{enumerate}
\end{lemma}

\begin{proof}
A chain of length zero has matrix \(I\), so consider a single chain
of positive length or a pair of chains. Expand the product by the
identifications of vertices across the two factors. Common edges
cancel. Let \(u,v,w\) be the left, joining and new right boundaries,
respectively. Of the set \(S\) of shared vertices, let \(U\) be those
incident to a remaining edge or equal to \(u\) or \(w\). The quotient
has \(v_P+v_\sigma-\abs S\) vertices, where \(v_P,v_\sigma\) are the
vertex counts of the two factors. Its isolated interior vertices
are precisely \(S\setminus U\), since the joining vertex \(v\) is
shared and all vertices of the factors outside the boundaries
have incident edges. The product normalization divides by
\(n^{(v_P+v_\sigma-2)/2}\) for one chain and
\(n^{(v_P+v_\sigma-1)/2}\) for two chains. Thus if \(\iota\) is the
boundary connectivity parameter of the quotient, the exponent of
\(n\) in \eqref{eq:rough} after normalization is
\[
 \frac{2-\abs U-\iota}{2}\quad\text{for one chain},\qquad
 \frac{1-\abs U-\iota}{2}\quad\text{for two chains}.
\]
When all edges survive, \(U=S\). For a single chain the boundaries
are connected through \(v\), so the only term with exponent zero is
the gluing with \(S=\{v\}\). For two chains this gluing also has
exponent zero, and an additional identification gives a negative
exponent.

When some but not all edges of \(\sigma\) cancel, its two edge sets
have at least two mixed vertices in \(U\). The only case to check is
one chain with \(\abs U=2\). The uncanceled edges of \(\sigma\) connect
the mixed vertices and \(w\). If all edges at \(w\) had canceled,
\(w\) would be a third vertex in \(U\). The left boundary \(u\) also
connects to \(U\). If it
is shared it belongs to \(U\), and otherwise a path in the old chain
to its first shared vertex reaches \(U\) through a surviving edge.
Thus \(\iota=1\) and the exponent is negative.

When all edges of \(\sigma\) cancel, they lie in the last block. Indeed,
a spanning cycle of \(\sigma\) cannot cross an articulation vertex of
the chain, and \(\sigma\) contains the right endpoint \(v\). The same
conclusion holds for a single edge. If the inclusion in the last block
is proper, apply the split to that block's edges. For \(\abs U=2\) the
remaining edges of the block join both mixed vertices, one of which is
\(w\). As above, in the single-chain case \(u\) connects to \(U\), so the
exponent is negative. If the cancellation is exact, deleting that
block leaves its left endpoint and the new right boundary among the
shared vertices. These coincide in the first case of rule 2. With two
chains and a right chain of length one, only the new right boundary
need be retained. These are precisely the exponent-zero cases.
The new right boundary must be a cycle neighbor of \(v\), or the
other vertex for a single edge. For each candidate, the two ordered
boundaries and the unique spanning cycle determine at most one
identification.

Removing \(z\) isolated interior vertices multiplies the unnormalized
matrix on the remaining vertices by \(n^z(1+O(n^{-1}))\), since their
labels are chosen without replacement. In the exponent-zero cases
this gives the claimed normalized terms. Their norms are at most
polylogarithmic by \eqref{eq:rough}, so the scalar errors and the terms
with negative exponents have norm \(o(1)\).
\end{proof}

\subsection{A positive semidefinite certificate}

Write \(P_j(x)=U_j(x/2)\) for the Chebyshev polynomials of the second
kind, orthonormal for the semicircle probability measure \(\mu\) on
\([-2,2]\). They satisfy \(P_j(2\cos t)=\sin((j+1)t)/\sin t\).
The connection between these polynomials and chain concatenation
comes from Potechin and Xu~\cite{PX}.
For a symmetric coefficient array \(c\) of square sum one,
let \(M_c=\sum_\alpha c_\alpha\widehat M_\alpha\), and write \(C_j(M_c)\)
for the sum of chains of \(j\) blocks with their product coefficients.
Here symmetry of the array means invariance under boundary reversal.
Set \(C_0(M_c)=I\). In the product \(C_j(M_c)M_c\), the appending term is
\(C_{j+1}(M_c)\) and the backtracking term is \(C_{j-1}(M_c)\), with
coefficient \(\sum_\alpha c_\alpha^2=1\). The remaining terms are
pairs of chains. Such pairs stay in the same class upon right
multiplication, up to the error in Lemma~\ref{lem:product}. The
recurrence \(P_{j+1}(x)=xP_j(x)-P_{j-1}(x)\), together with
\eqref{eq:sharp}, gives for each fixed \(j\)
\begin{equation}\label{eq:cheb}
 P_j(M_c)=C_j(M_c)+D_j(M_c)+o(1),
\end{equation}
where \(D_j(M_c)\) is a finite combination of two-chain matrices, and
the error is in operator norm with probability \(1-O(n^{-2})\).
In particular \(D_0=D_1=0\).

We need a nonnegative polynomial on \([-2,2]\) with a controlled
Chebyshev expansion. Let \(R_N(x)=\E\abs{4B/N-2}\) for a binomial
variable \(B\) with parameters \(N,(x+2)/4\). This is an even
polynomial and Jensen's inequality gives \(R_N(x)\ge\abs x\).
The binomial variance gives \(R_N(x)-\abs x\le2/\sqrt N\). The
derivative of the Bernstein polynomial is \(N/4\) times a binomial
average of adjacent differences of the data \(\abs{4j/N-2}\),
so \(\abs{R_N'}\le1\). Write
\[
 F_N(x)=x+R_N(x)=b_0+P_1(x)+\sum_{j=2}^N b_jP_j(x).
\]
The sine formula gives \((4-x^2)P_j''-3xP_j'+j(j+2)P_j=0\).

Integration by parts against \(\mu\) now yields
\(\sum_{j\ge2}j(j+2)b_j^2=\int(4-x^2)(R_N')^2\,d\mu\). In particular,
\begin{equation}\label{eq:poly-coeff}
 \sum_{j\ge2}j b_j^2\le\frac34,
 \qquad a^2:=1-\sum_{j\ge2}b_j^2=1-\int R_N^2\,d\mu+b_0^2.
\end{equation}
By \eqref{eq:poly-coeff},
\(\sum_{j\ge2}b_j^2\le\tfrac12\sum_{j\ge2}j b_j^2\le3/8\), so we take
the positive root \(a\ge\sqrt{5/8}\). Since \(R_N\) converges uniformly
to \(\abs x\), we have \(b_0\to\int\abs x\,d\mu>0\) and
\(\int R_N^2\,d\mu\to\int x^2\,d\mu=1\). It follows that \(a/b_0\to1\).

Fix \(N\). Give the single edge coefficient \(a\). For each chain of
\(j\ge2\) already constructed basic shapes, add its boundary edge
and give the resulting shape coefficient \(b_j\) times the product
of its constituent coefficients. The spanning paths in the blocks
concatenate to give the new cycle together with the added edge.
Every chord stays within a single block's interval on this cycle,
which makes the result a basic shape. Deleting its boundary edge
recovers the ordered chain, whose blocks are its
maximal two-connected subgraphs and bridges. Hence every constructed
shape has a unique coefficient. The construction preserves boundary
reversal. Give the edge depth zero and each new shape depth one plus
the largest depth of its constituents. If \(\kappa_h\) is the squared
coefficient norm through depth \(h\), then
\[
 \kappa_0=a^2,\qquad \kappa_{h+1}=a^2+\sum_{j\ge2} b_j^2 \kappa_h^j.
\]
This map fixes one and has Lipschitz constant at most \(3/4\) on
\([0,1]\) by \eqref{eq:poly-coeff}. Thus \(\kappa_h\uparrow1\), and the
normalized finite arrays converge in \(\ell_2\) to the full array.

Let \(M_c\) have the coefficients through depth \(h\), divided by
\(\sqrt{\kappa_h}\). For any chain of length at least two or pair of
chains, write \(\operatorname{cl}\) for adding its boundary edge with
the normalization held fixed. This multiplies the matrix entry
\((i,j)\) by \(A_{ij}\). The closure of a pair of chains has norm
\(o(1)\) by \eqref{eq:rough}, since it now has connected boundaries.
Use this operation term by term on \(C_j(M_c)\) and \(D_j(M_c)\). By
construction the coefficient array of \(\sum_{j\ge2}b_j
\operatorname{cl}(C_j(M_c))\) approaches that of
\(M_c-a A/\sqrt{n\kappa_h}\) in \(\ell_2\) as \(h\to\infty\). Indeed,
closing a chain is injective on the ordered chain types, and the
\(j\)-fold tensor power of the coefficient array is continuous in
\(\ell_2\) for fixed \(j\). Let \(e_h\to0\) denote this coefficient
difference in norm. Reversing a chain preserves its product
coefficient and reverses its closure, so the matrix difference is
symmetric. Lemma~\ref{lem:sharp} bounds its norm by \(3e_h\) with
probability \(1-O(n^{-2})\).

The matrix
\[
 W_0=b_0 I+\frac{a}{\sqrt{\kappa_h}}\frac{A}{\sqrt n}
       +\sum_{j\ge2} b_j (\mathrm{id}+\operatorname{cl})(C_j(M_c)+D_j(M_c))
\]
has diagonal \(b_0\) and nonedge entries \(-a/\sqrt{n\kappa_h}\). By
\eqref{eq:cheb} its symmetric part is within \(3e_h+o(1)\) of
\(F_N(M_c)\) in operator norm with probability \(1-O(n^{-2})\).
For any \(\eta>0\), choose \(h\) so that \(3e_h<\eta/3\). The fixed
polynomial \(F_N\) is nonnegative on \([-2,2]\), so it is at least
\(-\eta/3\) on \([-2-\delta,2+\delta]\) for some \(\delta>0\). This
interval contains the spectrum of \(M_c\) with the same probability
by \eqref{eq:sharp}. Adding \(\eta I\) to the symmetric part of
\(W_0\) therefore gives a positive semidefinite matrix for sufficiently
large \(n\). Apply \eqref{eq:witness}, using \(\kappa_h\le1\), to get
\(\vartheta(G_n)\le 1+(b_0+\eta)\sqrt n/a\) on this event. On the
complement use \(\vartheta(G_n)\le n\). Letting \(n\to\infty\), then
\(\eta\downarrow0\), then \(N\to\infty\), proves the independent-edge
upper bound.

\section{Random circulants}\label{sec:circulant}

\subsection{The smoothed Fourier program}

Write \(\mathcal D\) for the connection set and \(\mathcal N\) for its
complement in \(\Z_n\setminus\{0\}\). Define \(\varepsilon_0=0\),
\(\varepsilon_s=-1\) on \(\mathcal D\), and \(\varepsilon_s=1\) on
\(\mathcal N\). Thus \(\varepsilon_s=\xi_{[s]}\) for \(s\ne0\), where
\([s]=\{s,-s\}\) and the \(\xi_{[s]}\) are independent fair signs.
Let \(\mathcal F_{ks}=e^{-2\pi i ks/n}\). For
\(\mathcal S\in\{\mathcal D,\mathcal N\}\), set
\[
 P_{\mathcal S}=n^{-1}\mathcal F\operatorname{diag}(\one_{\mathcal S})\mathcal F^*,
 \qquad \Pi=n^{-1}\one\one^{\mathsf T},\qquad K=P_{\mathcal N}-P_{\mathcal D}.
\]
The operators \(P_{\mathcal D},P_{\mathcal N},\Pi\) are mutually orthogonal
projections that sum to \(I\). They preserve real even vectors, and
\(K^2=I-\Pi\).

Averaging \(X\) in \eqref{eq:theta} over cyclic translations gives a
circulant matrix with entries \(x_{i-j}/n\), where \(x\) is real and even.
Its eigenvalues are \(y=\mathcal F x/n\). Thus \(y\) is a real even
probability vector by positivity and the trace constraint. Since
\(x=\mathcal F^*y\), the edge constraints are \(P_{\mathcal D}y=0\).
Conversely, each real even probability vector in this kernel gives a
feasible matrix with entries \((\mathcal F^*y)_{i-j}/n\).
The objective is \(\sum_s x_s=ny_0\)~\cite{BBDFKK}.
For probability vectors in \(\ker P_{\mathcal D}\) and
\(\ker P_{\mathcal N}\) the inner product is \(1/n\), because only
their constant projections contribute.

Following \cite{CHHWY}, fix \(0<\tau\le1\), put \(w_0=1\) and
\(w_k=\tau/n\) for \(k\ne0\), and write \(\omega=\sum_k w_k\).
There exist positive real even vectors \(p,q\) in \(\ker P_{\mathcal D}\)
and \(\ker P_{\mathcal N}\), respectively, satisfying
\[
 p_kq_k=w_k,\qquad \sum_k p_k=\sum_k q_k=\sqrt{n\omega}.
\]
To obtain them, minimize
\((\sum_k p_k)^2/(2n)-\sum_k w_k\log p_k\) over positive real even
\(p\in\ker P_{\mathcal D}\). The domain contains \(\one\).
On a sublevel set the sum \(S=\sum_k p_k\) is bounded above, since the
objective is at least \(S^2/(2n)-\omega\log S\). The logarithmic terms then
bound every coordinate away from zero, so the minimum is attained
in the interior. With \(q_k=w_k/p_k\), the gradient \(\Pi p-q\) is
real and even. Its orthogonality to the real even part of
\(\ker P_{\mathcal D}\) therefore puts it in
\(\operatorname{ran}P_{\mathcal D}\). This gives
\(P_{\mathcal N}q=0\) and \(\Pi q=\Pi p\), so
\(\omega=\langle p,q\rangle=S^2/n\). As \(q/\sqrt{n\omega}\) is a feasible
probability vector for the complement, \(ny_0\le\sqrt{n\omega}/q_0\)
for every feasible \(y\). Using \(p_0q_0=1\) yields
\begin{equation}\label{eq:barrier-upper}
 \vartheta(G_n)/\sqrt n\le\sqrt\omega\,p_0.
\end{equation}

Set \(f=p-q\) and \(\Phi(v)_k=\sqrt{v_k^2+4w_k}\). Then
\(f=K(p+q)=K\Phi(f)\). For any real even mean-zero vector \(v\), let
\(u=\Phi(v)\), \(r=v-Ku\), \(\widetilde p=(u+v)/2\), and
\[
 \Delta=\sum_k w_k(\widetilde p_k/p_k+p_k/\widetilde p_k-2)\ge0.
\]
Since \((u_k+v_k)(u_k-v_k)=4w_k\), the two ratio terms sum to
\(\tfrac12(\langle u,p+q\rangle-\langle v,f\rangle)\). Using
\(K^2=I-\Pi\) and \(K(p+q)=f\) gives
\[
 \Delta=\sqrt\omega\,(\zeta-2\sqrt\omega)-\tfrac12\langle r,f\rangle,
 \qquad \zeta=n^{-1/2}\sum_k u_k.
\]
Also \(\zeta^2-4\omega=\norm{v}_2^2-\norm{Ku}_2^2=2\langle v,r\rangle-\norm{r}_2^2\).
Dividing by \(\zeta+2\sqrt\omega\) bounds \(\Delta\) by
\((\norm v_2+\tfrac12\norm f_2)\norm r_2\).
The term at zero in \(\Delta\) is at least \((\log(\widetilde p_0/p_0))^2\)
by \(2(\cosh x-1)\ge x^2\). Since
\(\log\widetilde p_0=\operatorname{arsinh}(v_0/2)\),
\eqref{eq:barrier-upper} gives
\begin{equation}\label{eq:residual}
 \log\bigl(\vartheta(G_n)/\sqrt n\bigr)
 \le \tfrac12\log\omega+\tfrac12\abs{v_0}+\sqrt\Delta,
 \qquad 0\le\Delta\le(\norm v_2+\tfrac12\norm f_2)\norm r_2.
\end{equation}

\subsection{A Fourier kernel bound}

We need a bound on \(\norm f_2\). The bound below has the same order
as the kernel estimate of Bandeira et al.~\cite{BBDFKK}, which uses
restricted isometry for uniformly sampled Fourier rows~\cite{HR}.
Here we take nets for a Rademacher process on the inversion classes.

\begin{lemma}\label{lem:kernel}
With probability \(1-O(e^{-c(\log n)^2})\), every real vector \(x\)
whose Fourier transform vanishes on \(\mathcal D\) or on \(\mathcal N\)
satisfies \(\norm x_2\le C n^{-1/2}(\log n)^2\norm x_1\).
\end{lemma}

\begin{proof}
Let \(\nu=cn/(\log n)^4\) for a small fixed \(c>0\), and let
\(\mathcal A=\{x\in\mathbb R^n:\norm x_2\le1,\ \norm x_1\le\sqrt{\nu}\}\).
Write \(d_{\mathcal F}(x,y)=\norm{\mathcal F(x-y)}_\infty\) and
\[
 \mathcal R(x)=\frac1n\sum_{a\ne0}\varepsilon_a\abs{(\mathcal F x)_a}^2.
\]
For \(x,y\in\mathcal A\), Parseval and Cauchy--Schwarz give
\(\abs{\mathcal R(x)-\mathcal R(y)}\le2d_{\mathcal F}(x,y)\). The sum of
squares of the independent sign coefficients of this difference is
at most \(8d_{\mathcal F}(x,y)^2/n\).
Indeed, each class has size at most two, and
\(\sum_a (\abs{(\mathcal F x)_a}+\abs{(\mathcal F y)_a})^2\le4n\).
Hoeffding's inequality gives the corresponding subgaussian tail.

For \(0<u\le\sqrt{\nu}\), the set \(\mathcal A\) has a \(d_{\mathcal F}\)-net inside
\(\mathcal A\) with logarithmic size at most \(C \nu(\log n)^2/u^2\).
For each \(x\), write it as the expectation of a random vector in
\(\{0,\pm\sqrt{\nu} e_1,\ldots,\pm\sqrt{\nu} e_n\}\) using its \(\ell_1\)
bound. An average of \(\lceil C \nu\log n/u^2\rceil\) independent
copies approximates \(x\) within \(u/2\) in \(d_{\mathcal F}\) with positive
probability by Hoeffding's inequality at the \(n\) complex Fourier
coordinates. Count the possible averages and choose a point of
\(\mathcal A\) from each \(u/2\)-ball that meets it.

Take nets at radii \(u_j=2^{-j}\sqrt{\nu}\), starting with \(\{0\}\) and
stopping when \(u_j\le1/32\). There are \(O(\log n)\) scales. Apply the
sign tail to pairs at consecutive scales at distance at most \(3u_j\).
For each such pair the probability of an increment exceeding
\(C_1\sqrt{\nu/n}\log n\) in absolute value is at most
\(2\exp(-c_1 C_1^2 \nu(\log n)^2/u_j^2)\).
The number of pairs is at most \(\exp(C_2\nu(\log n)^2/u_j^2)\).
Taking \(C_1\) large and summing over scales using \(\nu/u_j^2=4^j\)
bounds the failure probability by \(O(e^{-c'(\log n)^2})\).
Telescoping from zero and using the deterministic bound for the
last approximation gives
\(\sup_{x\in\mathcal A}\abs{\mathcal R(x)}\le C_3\sqrt{\nu/n}(\log n)^2+1/16\le1/2\)
for small enough \(c\). A unit vector in \(\mathcal A\) whose transform
vanishes on either set would instead have
\(\abs{\mathcal R(x)}=1-\abs{(\mathcal F x)_0}^2/n\ge1-\nu/n\). Rescaling proves
the estimate.
\end{proof}

For a real even vector, membership in \(\ker P_{\mathcal D}\) means
that its Fourier transform vanishes on \(\mathcal D\), and likewise
for \(\mathcal N\). The masses of \(p,q\) therefore give, on the event
of Lemma~\ref{lem:kernel},
\begin{equation}\label{eq:f-bound}
 \norm f_2=O_\tau((\log n)^2).
\end{equation}

\subsection{Tree fields over the cyclic group}

Consider unordered rooted trees whose internal nodes have at least
two children. For such a tree \(T\), write \(d(T)\), \(l(T)\), and
\(\aut(T)\) for its number of nodes, number of leaves, and number of
root-preserving automorphisms. Assign a residue to each leaf, and
let \(r_v\) be the sum of the leaf residues below node \(v\). Call an
assignment admissible when all \(r_v\) are nonzero and all classes
\([r_v]\) are distinct. Define the standardized tree field~\cite{CHHWY} by
\begin{equation}\label{eq:tree}
 Z_T(k)=\frac{n^{-l(T)/2}}{\sqrt{\aut(T)}}
       \sum_{\text{admissible assignments}}
             e^{-2\pi i k r_{\rm root}/n}\prod_{v\in V(T)}\varepsilon_{r_v}.
\end{equation}
Negating the leaf residues shows that the exponential may be replaced
by its cosine and that the field is real and even. The nonzero root
residue makes its spatial mean zero. For the one-node tree \(L\),
\(Z_L=\mathcal F\varepsilon/\sqrt n=\sqrt n K e_0\). Each \(Z_T(k)\)
is a homogeneous multilinear polynomial of degree \(d(T)\) in the
independent signs. If \(U=\sum_I u_I\prod_{b\in I}\xi_b\), define
\(\partial_b U=\sum_{I\ni b}u_I\prod_{a\in I\setminus\{b\}}\xi_a\).

A frequency \(k\) is regular if \(2k\ne0\) in \(\Z_n\). Thus the
exceptional frequencies are zero and, in even order, \(n/2\).

\begin{lemma}\label{lem:tree-moments}
For trees \(T,S\) with at most \(D\) nodes, let
\(B_D=\exp(CD\log(2D))\) for a suitable absolute \(C\). Uniformly in
frequency and in the modulus, the second moments of the coordinates
are at most \(B_D\), \(\E(\partial_b Z_T(k))^2\le B_D/n\), and
\begin{equation}\label{eq:deriv-var}
 \Var\!\left(\frac{1}{d(T)}\sum_b
           \partial_b Z_T(k)\,\partial_b Z_S(j)\right)\le B_D/n.
\end{equation}
For regular \(k,j\),
\begin{equation}\label{eq:cov}
 \E[Z_T(k)Z_S(j)]=\one_{\{T\simeq S,\ j=\pm k\}}+O(B_D/\sqrt n).
\end{equation}
\end{lemma}

\begin{proof}
For a set \(A\) of nodes in a tree with \(l\) leaves, let \(\rho(A)\)
be the rank of the descendant-leaf indicator rows at those nodes.
Subtract from each selected row the rows of its immediate selected
descendants, using the original rows for the subtraction. This is
a triangular change of rows with diagonal entries one. The resulting
rows have disjoint supports, and the nonempty supports give
independent rows. In particular, prescribed
residues of the nodes in \(A\) leave at most \(n^{l-\rho(A)}\) possible
leaf assignments. Each new nonzero row has a pivot of one, so this
argument works modulo every \(n\). For a nonempty proper subset \(A\),
\begin{equation}\label{eq:rank}
 \rho(A)+\rho(A^{\mathrm c})\ge l+1.
\end{equation}
The two spans together contain every leaf row. If their intersection
were zero, every minimally dependent set consisting of a parent and
all its children would lie on one side of the partition. This would
place the entire tree on one side.

A sign support (the set of classes used by an assignment) fixes all
leaf residues once its node roles and orientations are specified.
This costs at most \(d(T)!2^{d(T)}\) choices. Prescribing one class at
one node leaves at most \(2n^{l(T)-1}\) assignments. If \(N_I\) is the
number of assignments with support \(I\), the squared coefficient is
at most \(n^{-l(T)}N_I^2\). By sign orthogonality, the second moment
uses the sum over all \(I\) and the squared derivative the sum over
\(I\ni b\). Use \(N_I\le d(T)!2^{d(T)}\) and bound the corresponding
sums of \(N_I\) by \(n^{l(T)}\) and \(2d(T)n^{l(T)-1}\).
For \eqref{eq:deriv-var}, the nonconstant part comes from unequal
supports with a common class. Fix the matching of roles and
orientations on the overlap. In one of the two trees the overlap
occupies a proper subset \(A\) of its roles. With leaf counts \(l,l'\),
the number of assignment pairs is at most \(n^{l+l'-\rho(A)}\), by
choosing the other tree's leaves first. For any fixed symmetric
difference of the supports, specify the roles and orientations of
its classes. The first tree then has at most \(n^{l-\rho(A^{\mathrm c})}\)
choices, and its shared classes determine the rest of the second
tree's labels. Each summand has absolute weight at most
\(D n^{-(l+l')/2}\), counting the derivative choices. For a fixed
overlap pattern, the squared Walsh norm is at most the largest
absolute coefficient times the sum of absolute coefficients.
Multiplying the two counts and the squared weight gives the power
\(n^{l-\rho(A)-\rho(A^{\mathrm c})}\). The role allocations and
orientations contribute a factor \(\exp(O(D\log(2D)))\). There are at most
this many overlap patterns as well, so the \(L^2\) triangle
inequality and \eqref{eq:rank} give \eqref{eq:deriv-var}.

For \eqref{eq:cov}, the two supports must match completely. Different
node counts give zero. For unequal leaf counts, specify the smaller
leaf assignment first, which gives a bound \(B_D/\sqrt n\). Suppose
then that the trees have the same leaf count \(l\). Assignments using
the even-order singleton occur at most \(Dn^{l-1}\) times before
matching. Their contribution is \(O(B_D/n)\). For all other admissible
assignments a match specifies a unique signed node permutation.
A permutation that fails to preserve the target's parent-sum relations
as identities in the source leaves imposes a nonzero integer linear
congruence that is a signed sum of at most \(D+1\) indicator rows.
Fixing every leaf but one with nonzero coefficient \(a\) leaves
at most \(\gcd(\abs a,n)\le D+1\) choices, so at most
\((D+1)n^{l-1}\) assignments solve it. A permutation preserving all
the relations maps the real spaces of node sums bijectively, since
they have the same dimension. Its integer identities hold modulo
\(n\) too. Each source assignment determines the target assignment
from the target's matched leaf values. In its contribution we may
sum over unrestricted leaf assignments with error \(O(B_D/n)\).
Zero and singleton values and coincidences up to sign each impose a condition
with a coefficient of magnitude one, since distinct descendant-leaf
sets are nested or disjoint.

If the target root is matched to a proper source node, the two root
residues are independent uniform residues. The two leaf-sum rows have
a minor of determinant one from a leaf inside that subtree and a
leaf outside. The average cosine product is then zero.
A root-to-root match between trees with more than one node preserves
the leaves. The bijection of the real node-sum spaces preserves all
minimal dependencies, up to signs. In a minimal dependent set,
some subtracted row vanishes.
The relation between that row's original node and its immediate
selected descendants must use the entire set. If the set contains
the root, it therefore consists of the root and an antichain covering
its leaves. Such a set is minimally dependent, since the antichain
has disjoint nonempty leaf sets. The unique largest set of this form
consists of the root and all the leaves, so the signed permutation
preserves it. The dependence also forces a common sign at the root
and the leaves. The other node rows are sums of their descendant
leaves. The leaf permutation therefore preserves the corresponding
sets, their containment order, and the common sign. These matches are
exactly rooted tree isomorphisms with a common sign, including the
two trivial signed matches for one-node trees. For each of the
\(2\aut(T)\) matches when \(T\simeq S\), the cosine product averages to
\(1/2\) if \(j=\pm k\), and to zero otherwise, at regular frequencies.
The normalizations in \eqref{eq:tree} now give \eqref{eq:cov}.
\end{proof}

\subsection{Gaussian comparison}

For a sign polynomial \(U\) of degree at most \(D\), we use the bound
\begin{equation}\label{eq:hypercontractivity}
 \norm U_{L^{2m}}\le (2m-1)^{D/2}\norm U_{L^2}
 \qquad(m\text{ a positive integer}).
\end{equation}
Indeed, the average of \((a+b)^{2m}\) and \((a-b)^{2m}\) is at most
\((a^2+(2m-1)b^2)^m\) by comparing the even binomial coefficients.
Writing \(U=U_0+\xi U_1\) in one sign and taking the \(L^m\) triangle
inequality bounds \(\norm U_{L^{2m}}^2\) by
\(\norm{U_0}_{L^{2m}}^2+(2m-1)\norm{U_1}_{L^{2m}}^2\).
Iteration bounds this by the sum of squared sign coefficients with
weight \((2m-1)^j\) in degree \(j\), which proves
\eqref{eq:hypercontractivity}.

\begin{lemma}\label{lem:comparison}
At one regular frequency, or at two regular frequencies distinct up
to sign, take tree coordinates with at most \(D\) nodes. The Gaussian
counterparts are independent standard normals, one for each tree type
and each frequency. Let \(V\) be a real linear combination of these
coordinates with coefficient \(\ell_1\) norm at most \(C_0 4^D\), and
let \(G_V\) be its Gaussian counterpart, of variance at most \(C_0\).
Suppose \(\psi\) has three continuous derivatives with
\(\abs{\psi^{(j)}(x)}\le C_\psi(1+\abs x)^{2M+3}\) for \(0\le j\le3\).
Then, for a constant \(C\) depending only on \(C_0\),
\begin{equation}\label{eq:comparison}
 \abs{\E \psi(V)-\E \psi(G_V)}\le C_\psi E_{D,M}/\sqrt n,
 \qquad E_{D,M}=\exp\!\bigl(C[(D+1)(M+1)]^2\bigr).
\end{equation}
In particular, for a polynomial of total degree at most \(2M\), the
expectation error is at most its absolute coefficient sum times
\(E_{D,M}/\sqrt n\), after increasing \(C\). At arbitrary frequencies
the expected absolute value of such a polynomial is at most its
absolute coefficient sum times \(E_{D,M}\).
\end{lemma}

\begin{proof}
Write \(V=\sum_i a_iY_i\), where \(Y_i\) is a homogeneous coordinate of
degree \(d_i\), and let \(A_1=\sum_i\abs{a_i}\). Flip a uniformly chosen
sign, producing \(Y_i'\) and \(V'\). If there are \(m_0\) signs, the
conditional mean of \(Y_i'-Y_i\) is \(-\lambda_iY_i\), with
\(\lambda_i=2d_i/m_0\). Exchangeability gives for any test function \(g\)
\[
 \E[Y_i g(V)]=\frac{1}{2\lambda_i}
       \E[(Y_i'-Y_i)(g(V')-g(V))].
\]
Put \(\sigma^2=\E G_V^2\) and take \(G_V\) independent of the signs.
For \(V_t=\sqrt t\,V+\sqrt{1-t}\,G_V\), differentiate
\(\E \psi(V_t)\) and integrate by parts in \(G_V\). The derivative is
\(\E[V \psi'(V_t)]/(2\sqrt t)-\sigma^2\E \psi''(V_t)/2\).
Apply the flip identity with \(g(x)=\psi'(\sqrt t\,x+\sqrt{1-t}\,G_V)\),
conditioning on \(G_V\). Its linear Taylor term gives
\(\sqrt t\,\E[\Gamma \psi''(V_t)]\) in \(\E[V\psi'(V_t)]\), where
\[
 \Gamma=\sum_{i,j}\frac{a_i a_j}{d_i}
                  \sum_b\partial_bY_i\,\partial_bY_j.
\]
Orthogonality gives \(\E\Gamma=\E V^2\), and Lemma~\ref{lem:tree-moments}
gives \(\norm{\Gamma-\sigma^2}_{L^2}\le C A_1^2 B_D/\sqrt n\).

For the Taylor remainder set \(\delta_b=\sum_i\abs{a_i}\norm{\partial_bY_i}_{L^2}\).
Then \(\max_b \delta_b\le A_1\sqrt{B_D/n}\) and
\(\sum_b \delta_b^2\le A_1^2 D B_D\), the latter because
\(\sum_b\E(\partial_bY_i)^2=d_i\E Y_i^2\).
At a flip of sign \(b\), the remainder uses three sign differences
and the third derivative of \(\psi\) along the segment between
\(V_t\) and \(\sqrt t\,V'+\sqrt{1-t}\,G_V\).
Use H\"older with three \(L^6\) norms for the sign differences and
\(L^2\) for the derivative. By \eqref{eq:hypercontractivity}, the
absolute remainder in \(\E[V\psi'(V_t)]\) is bounded by
\[
 C t C_\psi 5^{3D/2}\Bigl(\sum_b \delta_b^3\Bigr)
     \bigl(1+2\norm V_{L^{2r}}+\norm{G_V}_{L^{2r}}\bigr)^r,
 \qquad r=2M+3.
\]
Here the probability of choosing \(b\) cancels the factor \(m_0\) in
the flip identity, and \(V'\) at each fixed flip has the law of \(V\).
Both this bound after division by \(2\sqrt t\), and the covariance
error involving \(\Gamma-\sigma^2\), are at most \(C_\psi E_{D,M}/\sqrt n\)
by \eqref{eq:hypercontractivity} and the Gaussian moment bound
\(\norm{G_V}_{L^{2r}}\le\sigma\sqrt{2r}\).

The moment factors in these estimates are bounded by powers of
\(1+(2r-1)^{D/2}A_1\sqrt{B_D}+\sigma\sqrt{2r}\), with exponents at
most \(r\). These and the remaining degree factors fit in \(E_{D,M}\).
Integrating over \(0<t<1\) gives \eqref{eq:comparison}. The endpoint
limits follow by dominated convergence, using the polynomial growth
bound for \(\psi\) and the moments of \(V\) and \(G_V\).

Each monomial of degree \(d\le 2M\) is a polarization of the \(d\)-th
power of linear combinations with coefficient \(\ell_1\) norm at
most one. Explicitly, average
\((\prod_{i=1}^d\eta_i)(\sum_{i=1}^d\eta_iY_i/d)^d\) over independent
fair signs \(\eta_i\), and multiply by \(d^d/d!\). The derivative
bounds for these powers and the polarization factor fit in the
same exponential envelope. The bound at arbitrary frequencies
follows by H\"older and \eqref{eq:hypercontractivity}.
\end{proof}

\subsection{The Gaussian equation}

We now restrict to rooted trees with even numbers of children at
internal nodes. Let \(\gamma_T\) be independent standard normals
indexed by these trees, and write \(h_j=\mathrm{He}_j/\sqrt{j!}\) for
the normalized probabilists' Hermite polynomials. Grafting a new
root above a multiset of children defines the isometry \(\mathcal J\)
from the span of the positive even-order Hermite products to the
span of the nonleaf coordinates \(\gamma_T\)~\cite{CHHWY}. In detail,
\[
 \mathcal J\!\left(\prod_T h_{m_T}(\gamma_T)\right)=\gamma_{T_*},
\]
where \(T_*\) has \(m_T\) children of type \(T\), and \(\sum_T m_T\) is
positive and even. These are orthonormal families indexed by the
same multisets, so \(\mathcal J\) extends to their closed span.

Write \(\varphi_\tau(x)=\sqrt{x^2+4\tau}\). We construct a centered
Gaussian linear combination \(H\) satisfying
\begin{equation}\label{eq:gauss}
 H=2\gamma_L+\mathcal J\bigl(\varphi_\tau(H)-\E\varphi_\tau(H)\bigr)
\end{equation}
with a tail bound for its squared coefficients by tree size.
Choose \(\sigma_\tau>0\) so that \(\E\varphi_\tau(\sigma_\tau\,g)=2\sqrt{1+\tau}\)
for a standard normal \(g\). The left side is continuous and strictly
increasing from \(2\sqrt\tau\) to infinity. Jensen's inequality for
the square root gives \(\sigma_\tau^2\ge4\), and Jensen's inequality for the
Euclidean norm of \((\sigma_\tau\,\abs g, 2\sqrt\tau)\) gives \(\sigma_\tau^2\le2\pi\).
Let \(\beta_j\) be the Hermite coefficients of
\(\varphi_\tau(\sigma_\tau\,g)\). Parseval gives
\(\sum_{j\ge2}\beta_j^2=\sigma_\tau^2-4\), and odd coefficients vanish.
Integration by parts in the Gaussian density gives
\begin{equation}\label{eq:hermite-energy}
 \sum_{j\ge1} j\beta_j^2=\sigma_\tau^2\,\E\!\left[\varphi_\tau'(\sigma_\tau\,g)^2\right]<\sigma_\tau^2.
\end{equation}
Indeed, the coefficients of the derivative in \(g\) are
\(\sqrt{j}\beta_j\) in degree \(j-1\). Repeating this argument shows
that for any integer \(p\ge1\), the squared Hermite tail at degrees
above \(M\) for \(\varphi_\tau(\sigma g)\) is \(O_{\tau,p}(M^{-p})\),
uniformly for \(0<\sigma\le\sqrt{2\pi}\). The \(p\)-th derivative has bounded
Gaussian \(L^2\) norm on this range, and the squared coefficient in
degree \(j\) is weighted by \(j(j-1)\cdots(j-p+1)\) in the derivative
norm. The bounded derivatives follow by differentiating
\(\varphi_\tau(x)=\sqrt{x^2+4\tau}\) with \(\tau>0\) fixed.

Define numbers \(u_T\) recursively by \(u_L=2/\sigma_\tau\) and
\[
 u_{T_*}=\frac{\beta_j}{\sigma_\tau}
      \sqrt{\frac{j!}{\prod_T m_T!}}\prod_T u_T^{m_T},
 \qquad j=\sum_T m_T\ge2.
\]
The squares \(u_T^2\) are the probabilities of the unordered trees
in a branching process having zero offspring with probability
\(4/\sigma_\tau^2\), and \(j\ge2\) offspring with probability \(\beta_j^2/\sigma_\tau^2\).
The multinomial factor counts orderings of the children. The mean
number of offspring is less than one by \eqref{eq:hermite-energy},
so the expected total population is finite and \(\sum_Tu_T^2=1\).
The offspring law has finite moments of every order by the Hermite
tail bound. The same is true of the total population. To see this,
truncate the process at height \(h\) and expand the \(p\)-th power of
one plus the sum of the child populations. Terms involving a single
child to power \(p\) have coefficient equal to the mean offspring
number times the previous height's \(p\)-th moment. The remaining
terms have uniformly bounded expectation by induction on \(p\),
using the \(p\)-th moment of the offspring number. This bounds the
\(p\)-th moments uniformly in \(h\), and monotone convergence gives
the assertion.

Set \(H=\sigma_\tau\sum_Tu_T\gamma_T\), with convergence in \(L^2\).
The Hermite generating function expands \(h_j(\sum_Tu_T\gamma_T)\)
in the products \(\prod_T h_{m_T}(\gamma_T)\) with coefficients
\(\sqrt{j!/\prod_Tm_T!}\prod_Tu_T^{m_T}\). To extend the finite identity,
normalize finite truncations of \(u\) to have norm one. The resulting
Gaussian sums converge in \(L^{2j}\) for each fixed \(j\), since the
variance of the difference tends to zero. Their Hermite polynomials
converge in \(L^2\), and each product coefficient converges to the
displayed value. Applying the identity to the Hermite expansion proves
\eqref{eq:gauss} by the recursion for \(u_T\). If \(H_D\) is the projection
of \(H\) onto trees of at most \(D\) nodes, the population moments give
\begin{equation}\label{eq:tree-tail}
 \E\abs{H-H_D}^2=O_{\tau,p}(D^{-p})\quad\text{for every integer }p\ge1,
 \qquad 4\le\E H_D^2\le2\pi.
\end{equation}
The lower variance bound follows from the leaf term.

\subsection{Transfer to the random Fourier program}

For a Gaussian linear combination \(Y\) with finite support, write
\(Y[n]\) for the vector obtained by replacing \(\gamma_T\) by \(Z_T\).
Write \(\|z\|_{n,2}^2=n^{-1}\sum_k\abs{z_k}^2\). From
Lemma~\ref{lem:tree-moments}, for \(Y\) supported on trees of at most
\(D\) nodes,
\begin{equation}\label{eq:linear-transfer}
 \E\norm{Y[n]}_{n,2}^2=(1+O(B_D/\sqrt n))\E Y^2,
 \qquad \E\abs{Y[n](0)}^2\le B_D\E Y^2,
\end{equation}
with \(C\) in \(B_D\) enlarged. In fact, there are at most \(4^D\) tree
types through size \(D\), by encoding plane representatives by their
traversals. Cauchy--Schwarz bounds the squared coefficient \(\ell_1\) norm
by \(4^D\E Y^2\). There are at most two exceptional coordinates.

Let \(P_M\) be the degree-\(M\) Hermite projection of \(\varphi_\tau\) under
the law of \(H_D\) in \eqref{eq:tree-tail}, and put
\(Y=\mathcal J(P_M(H_D)-\E P_M(H_D))\). With the constant in \(E_{D,M}\)
enlarged depending on \(\tau\), the estimates are
\begin{align}
 \E\norm{P_M(H_D[n])-\varphi_\tau(H_D[n])}_{n,2}^2
   &\le C_{\tau,p} M^{-p}+E_{D,M}/\sqrt n, \label{eq:activation}\\
 \E\norm{K P_M(H_D[n])-Y[n]}_{n,2}^2
   &\le E_{D,M}/\sqrt n. \label{eq:grafting}
\end{align}
For \eqref{eq:activation}, apply Lemma~\ref{lem:comparison} to
\(\psi(x)=(P_M(x)-\varphi_\tau(x))^2\). The coefficient bound for \(H_D\) in
\eqref{eq:linear-transfer} applies to the linear combination. The
coefficients of \(P_M\) in the normalized Hermite basis have squared
sum at most \(\E\varphi_\tau(H_D)^2\le2\pi+4\tau\). In the power basis
the absolute coefficient sum is at most
\(\exp(O_\tau(M\log(2M)))\), by the Hermite generating function and
\(\E H_D^2\ge4\). This bounds \(\psi\) and its first three derivatives as
required, up to a factor \((1+\abs x)^{2M+3}\). At the exceptional
frequencies the same growth bound and \eqref{eq:hypercontractivity}
give a total contribution at most \(E_{D,M}/n\) to the spatial
average. The Gaussian mean of \(\psi\) is the Hermite tail bounded above
in the proof of \eqref{eq:tree-tail}.

For \eqref{eq:grafting}, first take a normalized Hermite product
of positive even total degree at most \(M\) evaluated at the tree
fields, and denote the resulting vector by \(U\). With
\(\mathcal F_{ks}=e^{-2\pi i ks/n}\), the action of \(K\) multiplies the
mode \(k\mapsto e^{-2\pi i ks/n}\) by \(\varepsilon_s\). Expanding
\(K U\) therefore adds a sign at the new root, whose residue is the
sum of the child residues. The component of degree
\(1+\sum_T m_T d(T)\) is exactly \(Z_{T_*}\). Indeed, a zero root
residue gives zero, any repeated class reduces the degree by at
least two, and lower Hermite terms have smaller sign degree. The
leading term expands over all assignments to the individual copies
of the child trees, once per assignment to the grafted tree's
leaves. Its normalizing factor agrees because
\[
 \aut(T_*)=\prod_T m_T!\,\aut(T)^{m_T}
\]
for the multiplicities \(m_T\) of the child types. In particular, the
singleton class at even order is treated by the same sign rule.
By orthogonality of sign degrees and \(K^2=I-\Pi\),
\begin{align*}
 \E\norm{K U-Z_{T_*}}_{n,2}^2
   &= \E\norm{K U}_{n,2}^2-\E\norm{Z_{T_*}}_{n,2}^2,\\
 \E\norm{K U}_{n,2}^2
   &= \frac1n\sum_k\E U(k)^2-\frac1{n^2}\sum_{k,j}\E[U(k)U(j)].
\end{align*}
In the first sum the Gaussian expectation is one at regular
frequencies. In the double sum it is zero at regular frequencies
distinct up to sign. There are at most two excluded single terms
and \(O(n)\) excluded pairs. The absolute
power coefficient sum for a normalized Hermite product of total
degree at most \(M\) is \(\exp(O(M\log(2M)))\), again by the generating
function. Lemma~\ref{lem:comparison} bounds the average errors on
the regular terms by \(E_{D,M}/\sqrt n\) and the excluded
contributions by \(E_{D,M}/n\), with the constant enlarged.
The grafted tree has
at most \(1+DM\) nodes, so \eqref{eq:linear-transfer} puts its averaged
squared norm within the same error of one. This proves the bound
for each product.

The constant in the Hermite expansion of \(P_M(H_D)\) is annihilated by
\(K\). There are at most \((1+4^D+M)^M\) remaining products, and their
squared coefficient sum is at most \(2\pi+4\tau\). Cauchy--Schwarz
therefore gives \eqref{eq:grafting}, still within the envelope
\(E_{D,M}\) after increasing its constant.

Take \(v=H_D[n]/\sqrt n\). Its mean across the coordinates is zero,
it is real and even, and \eqref{eq:linear-transfer} bounds its second
moment and \(\E\abs{v_0}^2\). By \eqref{eq:gauss}, the isometry of
\(\mathcal J\), and the Lipschitz bound \(\abs{\varphi_\tau'}\le1\),
\[
 \E\abs{H_D-2\gamma_L-Y}^2=O_{\tau,p}(D^{-p}+M^{-p}).
\]
This linear combination is supported on trees of at most \(1+DM\)
nodes, so \eqref{eq:linear-transfer} transfers its bound to the spatial
mean square, with additional error at most \(E_{D,M}/\sqrt n\).
At zero the smoothing weight changes. In vector form
\[
 \Phi(v)=\varphi_\tau(H_D[n])/\sqrt n+(2+\beta)e_0,
 \qquad \abs\beta\le 2\abs{v_0}+2\sqrt{\tau/n}.
\]
The bound follows from \(\sqrt{v_0^2+4}-2\le\abs{v_0}\) and
\(\sqrt{v_0^2+4\tau/n}\le\abs{v_0}+2\sqrt{\tau/n}\).
Since \(Z_L=\sqrt n K e_0\), the residual is
\begin{align*}
 v-K\Phi(v)
 &=\frac{(H_D-2\gamma_L-Y)[n]}{\sqrt n}
   +\frac{Y[n]-K P_M(H_D[n])}{\sqrt n}\\
 &\quad+\frac{K\bigl(P_M(H_D[n])-\varphi_\tau(H_D[n])\bigr)}{\sqrt n}
   -\beta K e_0.
\end{align*}
The last term has expected squared norm \(O_\tau(B_D/n)\) by
\eqref{eq:linear-transfer} and \(\norm K_{\op}\le1\).
Bounding the other terms by the linear transfer,
\eqref{eq:grafting}, and \eqref{eq:activation} gives
\[
 \E\norm{v-K\Phi(v)}_2^2
 \le C_{\tau,p}(D^{-p}+M^{-p})+E_{D,M}/\sqrt n.
\]
Choose \(D=M=\lfloor(\log n)^{1/16}\rfloor\) and take \(p=200\).
Then \(E_{D,M}=n^{o(1)}\), and we have
\begin{equation}\label{eq:approx}
 \E\norm v_2^2=O_\tau(1),\qquad
 (\E\norm{v-K\Phi(v)}_2^2)^{1/2}=O_\tau((\log n)^{-6}),\qquad
 \E\abs{v_0}^2\le n^{-1+o(1)}.
\end{equation}

\subsection{Expectation bound}

Write \(\mathcal E_n\) for the event of Lemma~\ref{lem:kernel}.
By \eqref{eq:residual}, \eqref{eq:f-bound}, \eqref{eq:approx}, and
Cauchy--Schwarz, \(\E[\Delta\,\one_{\mathcal E_n}]=O_\tau((\log n)^{-4})\).
Set \(\delta_n=(\log n)^{-1/2}\). The event \(\mathcal E_n\), \(\abs{v_0}\le\delta_n\)
and \(\Delta\le\delta_n^2\) hold together with failure probability
\(O_\tau((\log n)^{-3})\), by Lemma~\ref{lem:kernel} and Markov's
inequality. On their intersection,
\eqref{eq:residual} gives
\(\vartheta(G_n)/\sqrt n\le\sqrt{1+\tau}\exp(3\delta_n/2)\).

For the exceptional contribution, the circulant matrix with entries
\(\varepsilon_{i-j}\), plus \(I\), agrees with \(J\) on the diagonal
and on nonedges. Its largest eigenvalue is at most
\(1+\norm{\mathcal F\varepsilon}_\infty\), which bounds
\(\vartheta(G_n)\) by \eqref{eq:theta}. Each Fourier coefficient is a
real sign sum with squared coefficient sum at most \(2n\), since
reflected pairs contribute coefficients of magnitude at most two
and the singleton contributes at most one. Hoeffding's inequality and a union
bound give \(\norm{\mathcal F\varepsilon}_\infty\le4\sqrt{n\log n}\)
with failure probability at most \(2n^{-3}\). Using this bound on the
exceptional event costs \(O_\tau((\log n)^{-5/2})\) in the normalized
expectation, and the spectral tail costs at most \(2n^{-5/2}\) by
\(\vartheta(G_n)\le n\). Hence
\(\limsup_{n\to\infty}\E\vartheta(G_n)/\sqrt n\le\sqrt{1+\tau}\).
Letting \(\tau\downarrow0\) and using the complement lower bound
completes the proof of Theorem~\ref{thm:main}.

\section*{Acknowledgments}

This work was carried out through automated proof discovery using
GPT-6.2 Sol Max in Codex. The human author sorted and reorganized the
resulting material and takes full responsibility for the work.
The \LaTeX{} manuscript and Lean formalization code are available at
\url{https://github.com/jinshanmu/LovaszTheta}.

\begin{thebibliography}{9}

\bibitem{Lovasz}
L. Lov\'asz, \emph{On the Shannon capacity of a graph},
IEEE Trans. Inform. Theory \textbf{25} (1979), no.~1, 1--7.
\doi{10.1109/TIT.1979.1055985}.

\bibitem{Juhasz}
F. Juh\'asz, \emph{The asymptotic behaviour of Lov\'asz' \(\vartheta\)
function for random graphs}, Combinatorica \textbf{2} (1982), no.~2,
153--155.
\doi{10.1007/BF02579314}.

\bibitem{PX}
A. Potechin and J. Xu, \emph{Sharp Lov\'asz-Theta Bounds on Random Graphs},
arXiv:2609.30064v1, 2026. \doi{10.48550/arXiv.2609.30064}.

\bibitem{BBDFKK}
A. S. Bandeira, J. B{\l}asiok, D. Dmitriev, U. Faure, A. Kireeva,
and D. Kunisky, \emph{The Lov\'asz number of random circulant graphs},
Proc. Int. Conf. Sampling Theory and Applications (SampTA), 2025,
1--5.
\doi{10.1109/SampTA64769.2025.11133544}.

\bibitem{CHHWY}
X. Chen, H. Hu, Y. Hu, W. Wang, and L. Yang, \emph{Sharp asymptotics for
the Lov\'asz number of random circulant graphs}, arXiv:2610.22971v1, 2026.
\doi{10.48550/arXiv.2610.22971}.

\bibitem{AMP}
K. Ahn, D. Medarametla, and A. Potechin, \emph{Graph Matrices: Norm
Bounds and Applications}, arXiv:1604.03423v5, 2021.
\doi{10.48550/arXiv.1604.03423}.

\bibitem{HR}
I. Haviv and O. Regev, \emph{The restricted isometry property of
subsampled Fourier matrices}, Proc. 27th Annu. ACM-SIAM Symp.
Discrete Algorithms (SODA), 2016, 288--297.
\doi{10.1137/1.9781611974331.ch22}.

\end{thebibliography}
\end{document}
